\section{习题}
\label{sec:jacobian-exercises}
% ============================================================================

这一章的真正难点，不是记住定义，而是学会在三个层面之间来回切换：
一会儿把 Jacobian 看成周期格的商，
一会儿把它看成曲线的次数零除子类群，
一会儿又要把 Hecke 特征值读成 Abel 簇上的几何分解。
下面五题正是按这条主线安排的：
前两题把“亏格 $1$ 曲线就是自己的 Jacobian”真正算出来；
第三题练习 Hecke 对应与极化的相容性；
第四题把不同特征形式对应的因子从 Jacobian 中拆开；
第五题则展示级别 $37$ 时 Jacobian 已经变成真正的二维 Abel 簇。

\begin{figure}[ht]
\centering
\begin{tikzpicture}[>=Stealth, font=\small, node distance=1.5cm]
  \node[draw, rounded corners, fill=blue!8, minimum width=4.2cm, minimum height=0.9cm] (a) {习题 1--2：亏格 $1$ 与 $J_0(11)$};
  \node[draw, rounded corners, fill=green!8, below=of a, minimum width=4.2cm, minimum height=0.9cm] (b) {习题 3：Hecke 对应与 Rosati};
  \node[draw, rounded corners, fill=orange!10, below=of b, minimum width=4.2cm, minimum height=0.9cm] (c) {习题 4：特征形式因子分解};
  \node[draw, rounded corners, fill=purple!10, below=of c, minimum width=4.2cm, minimum height=0.9cm] (d) {习题 5：二维 Jacobian $J_0(37)$};

  \draw[->, thick] (a) -- (b);
  \draw[->, thick] (b) -- (c);
  \draw[->, thick] (c) -- (d);
\end{tikzpicture}
\caption{本章习题路线图：先巩固“亏格 $1$ 曲线就是自己的 Jacobian”，再练习 Hecke 对应与极化，最后把不同特征形式对应的 Abel 因子从 $J_0(N)$ 中拆出来。}
\label{fig:ch06-exercises-roadmap}
\end{figure}


\begin{exercise}
\label{ex:jacobian-ec-self}
设 $E/\C$ 是椭圆曲线，零元为 $O$。证明 $E$ 是其自身的 Jacobian。
\end{exercise}

\begin{proof}[解答]
把 $E$ 看成一条亏格 $1$ 的紧黎曼面。
由\cref{thm:genus1-jacobian}，以 $O$ 为基点的 Abel--Jacobi 映射
\[
  \phi_O:E\longrightarrow J(E)
\]
是双全纯同构。
因此作为复环面，$E\cong J(E)$。

若想把这个同构写得更具体，可以这样看：
对任意点 $P\in E$，映射
\[
  P\longmapsto [P-O]
\]
把点送到次数为 $0$ 的除子类。
在 Abel 定理给出的识别 $J(E)\cong \operatorname{Pic}^0(E)$ 下，
这正是 Abel--Jacobi 映射。
它把单位元 $O$ 送到零类，并与曲线的群结构相容。
所以椭圆曲线不仅与自己的 Jacobian 双全纯同构，而且这个同构正是熟悉的“点减去基点”构造。
\end{proof}

\begin{exercise}
\label{ex:jacobian-x011}
计算 $g\bigl(X_0(11)\bigr)$，并验证 $J_0(11)$ 是一维 Abel 簇，也就是一条椭圆曲线。
\end{exercise}

\begin{proof}[解答]
由\cref{prop:gamma0N-ingredients} 中的 genus formula，
\[
  g\bigl(X_0(N)\bigr)=1+\frac{\mu}{12}-\frac{\nu_2}{4}-\frac{\nu_3}{3}-\frac{\nu_\infty}{2}.
\]
对 $N=11$，有
\[
  \mu=11\left(1+\frac1{11}\right)=12,
  \qquad \nu_2=0,
  \qquad \nu_3=0,
  \qquad \nu_\infty=2.
\]
代入得
\[
  g\bigl(X_0(11)\bigr)=1+\frac{12}{12}-0-0-\frac22=1.
\]

因此
\[
  \dim \Sk_2\bigl(\GammaO(11)\bigr)=g\bigl(X_0(11)\bigr)=1.
\]
再由\cref{prop:j0n-from-s2}，
\[
  \dim J_0(11)=\dim \Sk_2\bigl(\GammaO(11)\bigr)=1.
\]
所以 $J_0(11)$ 是一维 Abel 簇。
最后由\cref{prop:dim1-torus-abelian}，一维 Abel 簇就是椭圆曲线。
故 $J_0(11)$ 的确是一条椭圆曲线。
\end{proof}

\begin{exercise}
\label{ex:rosati-selfadjoint-hecke}
设 $J_0(N)$ 带有其自然主极化，记相应的 Rosati 对合为 $\dagger$。
证明由 Hecke 对应得到的算子满足
\[
  T_n^{\dagger}=T_n.
\]
\end{exercise}

\begin{proof}[解答]
把 $T_n$ 看成对应
\[
  X_0(N) \xleftarrow{\ \alpha_n\ } X_0(N;n) \xrightarrow{\ \beta_n\ } X_0(N)
\]
诱导的 Jacobian 自同态 $(\beta_n)_*\alpha_n^*$。
对主极化而言，Rosati 对合正好把一个对应换成它的转置对应，因此
\[
  T_n^{\dagger}=(\alpha_n)_*(\beta_n)^*.
\]
所以要证 $T_n^{\dagger}=T_n$，只需说明 Hecke 对应与其转置是同一个几何对象。

现在考察 $X_0(N;n)$ 上的点 $(E,C,D)$。
它在转置对应中应当对应于
\[
  \bigl(E/D,(C+D)/D,E[n]/D\bigr).
\]
更方便的表述是：从一个次数 $n$ 的循环同源
\[
  \varphi:E\to E/D
\]
可以取其对偶同源
\[
  \widehat{\varphi}:E/D\to E.
\]
对偶同源再次具有次数 $n$，并把目标端的相应循环子群送回原来的源端。
因此“先选一个阶 $n$ 子群再取商”和“在商曲线上选对偶核再回去”给出同一族数据。
换句话说，交换两边箭头的转置对应与原 Hecke 对应通过“取对偶同源”自然同构。

既然 Rosati 对合就是把对应换成转置，而这里转置对应与原对应相同，便得到
\[
  T_n^{\dagger}=T_n.
\]
这就是 Hecke 算子关于 Jacobian 自然极化的自伴性。
\end{proof}

\begin{exercise}
\label{ex:Af-Ag-factors}
设 $f,g\in \Sk_2(\GammaO(N))$ 是彼此正交的归一化 Hecke 特征形式。
证明 $A_f$ 和 $A_g$ 都同种于 $J_0(N)$ 的 Abel 子簇因子。
\end{exercise}

\begin{proof}[解答]
先看 $A_f$。
由\cref{prop:cotangent-of-Af}，$q_f^*H^0(A_f,\Omega^1)$ 可识别为
$\Sk_2(\GammaO(N))$ 中由 $f$ 的 Galois 共轭张成的 Hecke 稳定子空间，记为 $W_f$。
利用第\ref{ch:hecke}章的 Petersson 内积，取正交补 $W_f^{\perp}$。
因为 Hecke 算子彼此交换且关于 Petersson 内积自伴，
$W_f^{\perp}$ 也是 Hecke 稳定的。

设 $B_f\subset J_0(N)$ 是对应于 $W_f$ 的复子环面，
$C_f\subset J_0(N)$ 是对应于 $W_f^{\perp}$ 的复子环面。
由于 $W_f\oplus W_f^{\perp}=H^0(J_0(N),\Omega^1)$，
加法映射
\[
  B_f\times C_f\longrightarrow J_0(N)
\]
在原点的微分是同构，因此其核有限，故这是一个同种。
于是 $B_f$ 是 $J_0(N)$ 的一个 Abel 子簇因子。

再看商映射 $q_f:J_0(N)\twoheadrightarrow A_f$。
它在余切空间上的拉回把 $H^0(A_f,\Omega^1)$ 识别为 $W_f$，
因此 $q_f$ 限制在 $B_f$ 上的微分是同构。
于是
\[
  q_f|_{B_f}:B_f\longrightarrow A_f
\]
是满射且核有限，也就是一条同种。
故 $A_f\sim B_f$，从而 $A_f$ 同种于 $J_0(N)$ 的一个 Abel 子簇因子。

同理，$A_g$ 也同种于由 $g$ 的 Galois 共轭张成的 Hecke 子空间所对应的 Abel 子簇因子。
因为 $f$ 与 $g$ 正交，它们对应的 Hecke 子空间互相正交，故得到的是 Jacobian 中彼此独立的两个因子。
所以 $A_f$ 和 $A_g$ 都同种于 $J_0(N)$ 的 Abel 子簇因子。
\end{proof}

\begin{exercise}
\label{ex:jacobian-x037}
对 $\Gamma_0(37)$，验证
\[
  g\bigl(X_0(37)\bigr)=2,
\]
并说明为什么 $J_0(37)$ 是二维 Abel 簇。
\end{exercise}

\begin{proof}[解答]
仍用 genus formula。
对 $N=37$，有
\[
  \mu=37\left(1+\frac1{37}\right)=38,
  \qquad
  \nu_2=0,
  \qquad
  \nu_3=0,
  \qquad
  \nu_\infty=2.
\]
因此
\[
  g\bigl(X_0(37)\bigr)=1+\frac{38}{12}-0-0-\frac22=1+\frac{19}{6}-1=2.
\]
于是
\[
  \dim \Sk_2\bigl(\GammaO(37)\bigr)=g\bigl(X_0(37)\bigr)=2.
\]
由\cref{prop:j0n-from-s2} 得
\[
  \dim J_0(37)=\dim \Sk_2\bigl(\GammaO(37)\bigr)=2.
\]
而 Jacobian 总是 Abel 簇，所以 $J_0(37)$ 是二维 Abel 簇。
这说明到了级别 $37$，模曲线的 Jacobian 已经不再是一条椭圆曲线，而是一个真正的高维 Abel 簇。
\end{proof}
